\subsection*{第1章の索引用語}

\begin{description}
  \item[\term{電気信号}] 電圧や電流の変化によって情報を運ぶ物理的な信号です。
  \item[\term{負数}] 0より小さい数です。
  \item[\term{符号ビット}] 符号付き数の正負の解釈に関係するビットです。
  \item[\term{ASCII}] 英数字や制御文字へ番号を割り当てる文字コードです。
  \item[\term{浮動小数点}] 符号、指数、仮数に分けて実数を近似表現する方式です。
  \item[\term{IEEE 754}] 浮動小数点の形式と演算を定める国際規格です。
  \item[\term{MOSFET}] 電界によって電流を制御する代表的なトランジスタです。
  \item[\term{スイッチ}] 条件に応じて電気的な導通を切り替える素子の抽象化です。
  \item[\term{NOR}] ORの結果を反転する論理ゲートです。
  \item[\term{多ビット加算}] 桁上がりを隣の桁へ渡しながら複数桁を加える処理です。
  \item[\term{データバス}] CPU、メモリ、周辺装置の間で値を運ぶ信号線群です。
  \item[\term{メインメモリ}] 実行中の命令とデータを置く主記憶です。
  \item[\term{ノイマン型}] 命令とデータを同じ記憶へ置く計算機構成です。
  \item[\term{CPU}] Central Processing Unitの略で、命令を実行し機械を制御する処理装置です。
  \item[\term{命令レジスタ}] 現在デコードまたは実行する命令を保持するレジスタです。
\end{description}
